\subsection{Euler products for indefinite binary spaces}
\begin{theorem}\label{th:binaryeuler}
For every indefinite binary parent the ordinary coefficient is nonmultiplicative.
For a rationally nonsplit parent the proper coefficient is nonmultiplicative as well.
For a rationally split parent the proper coefficient is multiplicative precisely when
the positive integer \(B\) of \cite[Thm.~1.1]{Hyperbolic} divides \(24\).
In each multiplicative case
\[
 Z_{SO}(s)=\zeta(s-1)\sum_{h\mid B}h^{-s}.
\]
Here multiplicativity is equivalent to the normalized Euler product of the prime-power subseries.
\end{theorem}
\begin{proof}
First assume the space is nonsplit, and let \(K=\mathbb Q(\sqrt{-\det B_L})\).
Let \(z,z^{-1}\) be the eigenvalues of the generator of the integral rotations
modulo scalar signs in Theorem~\ref{th:inertreal}. They are quadratic integral units
in \(K\), by the diagonal action on the two isotropic lines over \(K\).
Galois conjugation interchanges them. Put \(\alpha=z^2\).
At every good inert odd \(p\) outside the finitely many primes dividing
the discriminant of the characteristic polynomial of the generator,
the two reduced eigenvalues are distinct in \(\mathbb F_{p^2}\).
Consequently the projective order \(o_p=o_p(1)\) is the multiplicative order
of \(\alpha\) in that residue field. Choose such a \(p\) with \(o_p>1\);
infinitely many inert primes are available, whereas \(z^2-1\) is a nonzero
algebraic integer.

Choose a prime \(\ell\mid o_p\), let \(a=v_\ell(p^2-1)\), and put \(N=\ell^a\).
We obtain infinitely many further good inert primes \(q\) with
\(\ell\mid o_q\) by Chebotarev.
In the cyclic group \(\mathbb F_{p^2}^{\times}\) an element is an \(N\)th power
if and only if its order is prime to \(\ell\). Thus \(X^N-\alpha\) has no
root over \(\mathbb F_{p^2}\). Consider the finite extension
\[
 F=K(\zeta_{\ell N},\sqrt[N]{\alpha}).
\]
It is Galois over \(\mathbb Q\): it contains all the roots of \(X^N-\alpha\)
and the inverses of those roots give all the roots of \(X^N-\alpha^{-1}\).
The prime \(p\) is unramified in \(F\): it is unramified in
\(K\), is prime to \(\ell N\), and \(N\) and \(\alpha\)
are local units. The roots consequently generate an unramified
algebra by separable reduction.
Choose a prime of \(F\) above \(p\) with arithmetic Frobenius \(\sigma\).
Its restriction to \(K\) is nontrivial. Since the roots have distinct
reductions, \(\sigma^2\) has no fixed point on the roots of \(X^N-\alpha\).

The Chebotarev density theorem \cite{StevenhagenLenstra} gives
infinitely many rational primes \(q\) for which one can choose a
prime of \(F\) with Frobenius equal to \(\sigma\).
Exclude \(p\) and the finitely many exceptional primes.
The action on \(K\) makes \(q\) inert, and its action on \(\zeta_{\ell N}\)
gives \(q\equiv p\bmod\ell N\), hence \(v_\ell(q^2-1)=a\).
The same permutation \(\sigma^2\) on the roots shows that \(X^N-\alpha\)
has no root over \(\mathbb F_{q^2}\). Cyclicity again yields \(\ell\mid o_q\).

Now \(g=\gcd(o_p,o_q)\ge2\). Write \(A=(p+1)/o_p\) and \(C=(q+1)/o_q\).
The proper coefficients at \(p,q,pq\) are \(A,C,gAC\), by Theorem~\ref{th:inertreal}.
If the parent has no improper integral isometry these are also the ordinary coefficients.
Otherwise the ordinary coefficients are \((A+2)/2,(C+2)/2,(gAC+4)/2\).
Both \(A,C\) are at least two: at an odd inert prime \(z\) reduces to the norm-one
group of order \(p+1\), and \(z^2\) has order dividing \((p+1)/2\).
Four times the ordinary discrepancy at \(pq\) is
\[
 (2g-2)AC+(A-2)(C-2)>0 .
\]
Neither coefficient sequence is therefore multiplicative in the nonsplit case.

In the split plane every integral stabilizer is finite by
Proposition~\ref{prop:splitbinary}. The prime-extension limit of
Lemma~\ref{lem:pr} and its proof use the same local restoration and
finite stabilizers, so they apply here as well.
If the ordinary coefficient were multiplicative, the two large-prime
limits in the proof of Theorem~\ref{th:nonmult} would force
\(|\Gamma_O(L)|/2=1\) and then \(|\Gamma_O(M)|/2=1\) for every represented \(M\).
But a rational orthogonal basis gives a diagonal sublattice inside \(L\),
and changing one coordinate sign is a nonscalar automorphism of that lattice.
The resulting nonscalar automorphism contradicts ordinary
multiplicativity.

For the split proper coefficient use the normalization of
\cite[Thm.~1.1]{Hyperbolic}: embed a scalar copy of \(L\) into the plane
\(H=\mathbb Z^2\) with quadratic form \(xy\), and define
\(B\mathbb Z=[H:L](\mathfrak nL)^{-1}\), where \(\mathfrak nL\) is the ideal
generated by the integral values of \(xy\) on \(L\).
For \(b\mid B\) put
\[
 C_b(m)=\#\{d^2\bmod b:d\mid m,\ (d,b)=1\}.
\]
Their formula is
\[
 Z_{SO}(s)=\frac{\zeta(s-1)}{\zeta(s)}
     \sum_{b\mid B}(B/b)^{-s}\sum_{m\ge1}C_b(m)m^{-s}.
\]
If the unit group modulo \(B\) has exponent dividing two, all units
modulo every divisor \(b\) have square one, by surjectivity of reduction
on unit groups. Then \(C_b(m)=1\) for every \(b,m\); the asserted formula
follows, and its coefficients are multiplicative.
Otherwise take a unit \(v\bmod B\) with \(v^2\ne1\), and choose
distinct primes \(p,q\equiv v\bmod B\). Among indices coprime to \(B\)
only the undilated term \(b=B\) contributes. Thus on \(1,p,q,pq\)
the coefficients are the Dirichlet convolution of Euler's totient with \(C_B\).
Here \(C_B(p)=C_B(q)=2\), and \(C_B(pq)\le3\) since the only residues
are \(1,v^2,v^4\). Multiplicativity of the totient gives
\(a_{SO}(pq)-a_{SO}(p)a_{SO}(q)=C_B(pq)-4<0\).
The unit group modulo \(B\) has exponent at most two exactly when \(B\mid24\).
For an odd prime power \(r^e\mid B\), reduction modulo \(r\) forces \(r=3\);
the unit \(1+3\) rules out \(e\ge2\). At two every odd square is
one modulo eight and \(3^2\ne1\bmod16\). The Chinese remainder theorem combines these local restrictions.
\end{proof}
